\documentclass[12pt, a4paper]{article}
\usepackage{../../../myconfig} % Gọi file cấu hình từ ngoài gốc vào

% ==========================================
% BẮT ĐẦU NỘI DUNG TÀI LIỆU
% ==========================================
\begin{document}

% Truyền 3 tham số: {Nhãn cam} {Chữ to chính giữa} {Chữ ở góc phải Header}
\titlebox{Chủ đề: Hàm số}{Hàm logarit, mũ}{Hàm số: Logarit, mũ}

% --- CÂU 1 ---
\questionLinesManual{0}{1}{Rút gọn biểu thức \( P = x^{\frac{1}{3}} \sqrt[6]{x} \) với \( x > 0 \).}{
    \begin{tasks}(4)
        \task \( P = \sqrt{x} \)
        \task \( P = x^{\frac{1}{8}} \)
        \task \( P = x^{\frac{2}{9}} \)
        \task \( P = x^2 \)
    \end{tasks}
}

% --- CÂU 2 ---
\questionLinesManual{0}{2}{Cho \( a \) là số thực dương khác 1. Mệnh đề nào dưới đây đúng với mọi số dương \( x, y \)?}{
    \begin{tasks}(2)
        \task \( \log_a \dfrac{x}{y} = \log_a x - \log_a y \)
        \task \( \log_a \dfrac{x}{y} = \log_a (x - y) \)
        \task \( \log_a \dfrac{x}{y} = \log_a x + \log_a y \)
        \task \( \log_a \dfrac{x}{y} = \dfrac{\log_a x}{\log_a y} \)
    \end{tasks}
}

% --- CÂU 3 ---
\questionLinesManual{0}{3}{Với \( a, b \) là các số thực dương tùy ý và \( a \neq 1 \), \( \log_a \left( \sqrt[5]{b} \right) \) bằng}{
    \begin{tasks}(4)
        \task \( 5 \log_a b \)
        \task \( \dfrac{1}{5} + \log_a b \)
        \task \( 5 + \log_a b \)
        \task \( \dfrac{1}{5} \log_a b \)
    \end{tasks}
}

% --- CÂU 4 ---
\questionLinesManual{0}{4}{Nghiệm của phương trình \( 3^{x-1} = 27 \) là}{
    \begin{tasks}(4)
        \task $x = 4$
        \task $x = 3$
        \task $x = 2$
        \task $x = 1$
    \end{tasks}
}

% --- CÂU 5 ---
\questionLinesManual{0}{5}{Nghiệm của phương trình \( \log_3(2x-1) = 2 \) là}{
    \begin{tasks}(4)
        \task $x = 3$
        \task $x = 5$
        \task $x = \dfrac{9}{2}$
        \task $x = \dfrac{7}{2}$
    \end{tasks}
}

% --- CÂU 6 ---
\questionLinesManual{0}{6}{Tập nghiệm của bất phương trình \( \log x \geq 1 \) là}{
    \begin{tasks}(4)
        \task \( (10; +\infty) \)
        \task \( (0; +\infty) \)
        \task \( [10; +\infty) \)
        \task \( (-\infty; 10) \)
    \end{tasks}
}

% --- CÂU 7 ---
\questionLinesManual{0}{7}{Tập nghiệm của bất phương trình \( \left( \dfrac{1}{3} \right)^x > 9 \) là}{
    \begin{tasks}(4)
        \task \( (2; +\infty) \)
        \task \( (-\infty; -2) \)
        \task \( (-\infty; 2) \)
        \task \( (-2; +\infty) \)
    \end{tasks}
}

% --- CÂU 8 ---
\questionLinesManual{0}{8}{Tập xác định của hàm số \( y = \log_9(x - 3) \) là}{
    \begin{tasks}(4)
        \task \( (-\infty; 3) \)
        \task \( (3; +\infty) \)
        \task \( [3; +\infty) \)
        \task \( (-\infty; +\infty) \)
    \end{tasks}
}

% --- CÂU 9 ---
\questionLinesManual{0}{9}{Nghiệm của phương trình \( 5^x = 10 \) là:}{
    \begin{tasks}(4)
        \task $2$
        \task \( \log_5 10 \)
        \task $5$
        \task \( \ln 10 \)
    \end{tasks}
}

% --- CÂU 10 ---
\questionLinesManual{0}{10}{Tập nghiệm của bất phương trình \( \log(x-1) < 2 \) là}{
    \begin{tasks}(4)
        \task \( (-\infty; 101) \)
        \task \( (1; 101) \)
        \task \( (1; 3) \)
        \task \( (101; +\infty) \)
    \end{tasks}
}

\end{document}